\section{Addition of the intrinsic trace data}\label{sec:trace-additivity}

Return to a submanifold $Y$ of arbitrary dimension. Subadditivity of multiplier ideals gave one inequality for the trace of a sum. For the reverse inequality, first fix an analytic approximation to the sum, then compare it with finer approximations to the two summands using \cref{lem:analytic-minorant-product-ideals}.

\Needspace{8\baselineskip}
\begin{theorem}\label{thm:trace-data-additivity}
Let $u,v\in\QPSH(X)$ have zero generic Lelong number along the connected compact submanifold $Y$. Then, on every divisor over $Y$,
\[
 \tau_{Y,u+v}=\tau_{Y,u}+\tau_{Y,v}.
\]
\end{theorem}
\begin{proof}
The inequality $\geq$ was proved in \cref{prop:trace-profile-calculus}. Fix a sufficiently large integer $k$ and let $S_k$ be the global analytic approximant to $u+v$, of type $(1/k,\mathcal I(k(u+v)))$. By \cref{thm:demailly-hermitian-input}, $u+v\leq S_k$. Apply \cref{lem:analytic-minorant-product-ideals} with $s=S_k$ to the global analytic approximants $U_m,V_m$ of $u,v$, respectively, with types $(1/m,\mathcal I(mu))$ and $(1/m,\mathcal I(mv))$. Their restrictions and that of $S_k$ are genuine because all three original potentials have zero generic Lelong number along $Y$.

Restrict the resulting ordinary inequality to $Y$, pull back to any fixed smooth compact modification, and take the generic Lelong number along its divisor $E$. This gives
\[
 \frac{\operatorname{ord}_E\mathfrak a_m(u;Y)}m
 +\frac{\operatorname{ord}_E\mathfrak a_m(v;Y)}m
 \geq(1-\delta_m)
       \frac{\operatorname{ord}_E\mathfrak a_k(u+v;Y)}k.
\]
Let $m\to\infty$. Finiteness and the limit formula in \cref{thm:intrinsic-trace-profile} yield
\[
 \tau_{Y,u}(E)+\tau_{Y,v}(E)
 \geq\frac{\operatorname{ord}_E\mathfrak a_k(u+v;Y)}k.
\]
Now let $k\to\infty$. This proves the opposite inequality and hence equality.
\end{proof}

When $Y$ is a curve, the theorem gives
$D_{u+v}=D_u+D_v$. With one fixed nonpositive Green kernel the corresponding Green potentials add, by integration of positive measures. The atomic trace therefore satisfies
\[
 \operatorname{Tr}^{\mathrm{at}}_Y(u+v)
 \sim\operatorname{Tr}^{\mathrm{at}}_Y(u)
       +\operatorname{Tr}^{\mathrm{at}}_Y(v).
\]
Positive real homogeneity follows from homogeneity of the trace data and linearity of the Green operator.
